\documentclass[11pt,a4paper]{article}
\usepackage[utf8]{inputenc}
\usepackage[T1]{fontenc}
\usepackage[french]{babel}
\frenchbsetup{StandardLists=true}
\usepackage[left=2cm,right=2cm,top=2cm,bottom=2cm]{geometry}
\usepackage[dvipsnames]{xcolor}
\usepackage{fouriernc}
\usepackage{amsmath,amssymb,amsfonts}
\usepackage{array,tabularx,booktabs,multirow}
\usepackage{colortbl}
\usepackage{graphicx}
\usepackage{mdframed}
\usepackage{enumitem}
\usepackage{fancyhdr}
\usepackage{chemformula}
\usepackage{awesomebox}
\setlength{\aweboxvskip}{0mm}
\usepackage{tcolorbox}
\tcbuselibrary{skins}
\usepackage{fontawesome5}

\definecolor{exoheader}{HTML}{1E5AA0}
\definecolor{exolight}{HTML}{DCE8FF}
\definecolor{warnorange}{RGB}{220,100,0}
\definecolor{problemebg}{HTML}{FFFDE7}

\renewcommand{\arraystretch}{1.8}
\setlength{\extrarowheight}{1mm}

\pagestyle{fancy}
\renewcommand\headrulewidth{1pt}
\fancyhead[L]{Académie de Besançon}
\fancyhead[C]{\textbf{TP --- Synthèse d'un biocarburant}}
\fancyhead[R]{Terminale / BTS Chimie}
\fancyfoot[C]{page \thepage}
\fancyfoot[R]{2025 -- 2026}
\fancyfoot[L]{Alexis RUIZ}

\newcommand{\sectitle}[1]{%
  \vspace{6pt}
  \noindent\colorbox{exoheader}{\parbox{\dimexpr\linewidth-2\fboxsep}%
    {\color{white}\bfseries\large\strut\quad #1}}
  \vspace{4pt}
}
\newcommand{\subsectitle}[1]{%
  \vspace{4pt}
  \noindent\colorbox{exolight}{\parbox{\dimexpr\linewidth-2\fboxsep}%
    {\color{exoheader}\bfseries\strut\quad #1}}
  \vspace{2pt}
}
\newcounter{question}
\newcommand{\question}{%
  \stepcounter{question}%
  \noindent\textcolor{exoheader}{\textbf{Question \thequestion.}}\quad
}

\begin{document}

\noindent\textbf{NOM :}\hfill\textbf{Prénom :}\hfill\textbf{Classe :}\hfill\phantom{x}\\[0.2cm]
\hrule\vspace{0.2cm}

\begin{center}
  {\large\bfseries\textcolor{exoheader}{Travaux Pratiques de Chimie Organique}}\\[4pt]
  {\LARGE\bfseries Synthèse d'un biocarburant (biodiesel)}\\[2pt]
  {\small Académie de Besançon --- Belfort}
\end{center}
\hrule\vspace{6pt}

\begin{tcolorbox}[enhanced, colback=problemebg, colframe=exoheader, boxrule=1.5pt,
  title={\faFlask\quad Problématique},
  fonttitle=\bfseries\color{white},
  attach boxed title to top left={yshift=-2mm, xshift=4mm},
  boxed title style={colback=exoheader, colframe=exoheader}]
\textit{Face à la raréfaction des ressources fossiles, les biocarburants représentent
une alternative. Le biodiesel est produit par \textbf{transestérification} d'huile végétale avec un alcool.
\textbf{Comment synthétiser du biodiesel au laboratoire à partir d'huile de tournesol ?}}
\end{tcolorbox}

\vspace{6pt}

% ══════════════════════════════════════════════════════════════════════════════
\sectitle{Partie I --- Étude de la réaction}
% ══════════════════════════════════════════════════════════════════════════════

\begin{mdframed}[backgroundcolor=exolight, linecolor=exoheader, linewidth=1.2pt]
La \textbf{transestérification} est la réaction d'un triglycéride (ester d'acides gras)
avec un alcool (éthanol ou méthanol) en présence d'un catalyseur basique (\ch{NaOH}).
Elle produit un \textbf{ester éthylique} (biodiesel) et du \textbf{glycérol}.
\end{mdframed}

\vspace{4pt}
\question L'huile de tournesol est considérée comme du trioléate de glycérol.
Écrire l'équation de transestérification avec l'éthanol. \dotfill

\vspace{0.3cm}\noindent\dotfill

\vspace{0.3cm}\noindent\dotfill

\question Pourquoi utilise-t-on un excès d'éthanol lors de la synthèse ? \dotfill

\vspace{0.3cm}\noindent\dotfill

\question Les deux biodiesel EEHV (ester éthylique) et EMHV (ester méthylique) ont des
performances similaires. Quel critère orientera le choix du type de biocarburant utilisé
dans un moteur Diesel ? (utiliser la notion d'indice de cétane) \dotfill

\vspace{0.3cm}\noindent\dotfill

% ══════════════════════════════════════════════════════════════════════════════
\sectitle{Partie II --- Protocole expérimental}
% ══════════════════════════════════════════════════════════════════════════════

\begin{mdframed}[backgroundcolor=orange!15, linecolor=warnorange, linewidth=2pt]
\textbf{\textcolor{warnorange}{\faExclamationTriangle\quad ATTENTION}} ---
\ch{NaOH} est corrosif. L'éthanol est inflammable. Porter \textbf{lunettes et gants}.
Pas de flamme nue lors des manipulations.
\end{mdframed}

\vspace{4pt}
\subsectitle{Étape 1 --- Synthèse du biodiesel}

Conditions opératoires : rapport molaire alcool/huile = 6 ; catalyseur \ch{NaOH} à 0,7\% en masse ;
température 70°C ; agitation forte ; durée 45 min.

\vspace{4pt}
\question Rédiger un protocole expérimental complet pour synthétiser le biodiesel
à partir de 70~mL d'huile de tournesol (densité 0,92). \dotfill

\vspace{0.3cm}\noindent\dotfill

\vspace{0.3cm}\noindent\dotfill

\vspace{0.3cm}\noindent\dotfill

\question Quel est l'aspect du milieu réactionnel en début de synthèse ?
Pourquoi est-il nécessaire d'agiter fortement ? \dotfill

\vspace{0.3cm}\noindent\dotfill

\subsectitle{Étape 2 --- Décantation et séparation}

\begin{enumerate}[leftmargin=1.5cm, label=\textcolor{exoheader}{\textbf{\arabic*.}}]
  \item Après refroidissement, transférer le mélange dans une ampoule à décanter.
  \item Attendre la séparation en deux phases (laisser décanter 10 à 15 min).
  \item Identifier et nommer les deux phases.
  \item Récupérer la phase supérieure (biodiesel brut) dans un bécher.
\end{enumerate}

\question Identifier les deux phases dans l'ampoule à décanter.
Laquelle est la plus dense ? Justifier.

\vspace{0.3cm}
\begin{center}
\begin{tabularx}{0.7\linewidth}{|>{\centering\arraybackslash\columncolor{exolight}}X
                                |>{\centering\arraybackslash}X|}
\hline
\rowcolor{exoheader}
\color{white}\textbf{Phase} & \color{white}\textbf{Contenu} \\
\hline
Phase supérieure & \dotfill \\
\hline
Phase inférieure & \dotfill \\
\hline
\end{tabularx}
\end{center}

% ══════════════════════════════════════════════════════════════════════════════
\newpage
\sectitle{Partie III --- Valorisation du glycérol}
% ══════════════════════════════════════════════════════════════════════════════

\begin{mdframed}[backgroundcolor=exolight, linecolor=exoheader, linewidth=1.2pt]
Le \textbf{glycérol} (\ch{C3H8O3}) est un co-produit de la transestérification.
Sa valorisation est essentielle pour l'économie du procédé.
\end{mdframed}

\vspace{4pt}
\question En utilisant la formule semi-développée du glycérol, identifier les groupes
fonctionnels et préciser le nom de ce composé selon sa famille chimique. \dotfill

\vspace{0.3cm}\noindent\dotfill

\question Citer deux utilisations industrielles du glycérol. \dotfill

\vspace{0.3cm}\noindent\dotfill

% ══════════════════════════════════════════════════════════════════════════════
\sectitle{Partie IV --- Caractéristiques du biodiesel}
% ══════════════════════════════════════════════════════════════════════════════

\question Écrire l'équation de combustion complète de l'ester éthylique d'huile de tournesol
(EEHV, formule brute \ch{C20H38O2}). \dotfill

\vspace{0.3cm}\noindent\dotfill

\question Comparer les avantages et inconvénients du biodiesel par rapport au gazole fossile.

\vspace{4pt}
\begin{center}
\begin{tabularx}{\linewidth}{|>{\centering\bfseries\columncolor{exoheader}\color{white}}m{3cm}
                              |>{\centering\arraybackslash}X
                              |>{\centering\arraybackslash}X|}
\hline
\rowcolor{exoheader}
 & \color{white}\textbf{Avantages} & \color{white}\textbf{Inconvénients} \\
\hline
Biodiesel & & \\[1.5cm]
\hline
Gazole fossile & & \\[1.5cm]
\hline
\end{tabularx}
\end{center}

% ══════════════════════════════════════════════════════════════════════════════
\newpage
\sectitle{Annexe --- Données}
% ══════════════════════════════════════════════════════════════════════════════

\begin{center}
\begin{tabularx}{\linewidth}{|>{\centering\bfseries\color{white}\columncolor{exoheader}}m{3cm}
                              |>{\centering\arraybackslash}m{3cm}
                              |>{\centering\arraybackslash}m{2cm}
                              |>{\centering\arraybackslash}X|}
\hline
\rowcolor{exoheader}
\color{white}\bfseries Espèce & \color{white}\bfseries Formule brute &
\color{white}\bfseries $M$ (g/mol) & \color{white}\bfseries Densité / remarques \\
\hline
Trioléate de glycérol & \ch{C57H104O6} & 884 & $d = 0{,}92$ \\
\hline
Éthanol & \ch{C2H5OH} & 46 & $d = 0{,}789$ ; inflammable \\
\hline
EEHV (biodiesel) & \ch{C20H38O2} & 310 & $d \approx 0{,}88$ \\
\hline
Glycérol & \ch{C3H8O3} & 92 & $d = 1{,}26$ \\
\hline
\ch{NaOH} & \ch{NaOH} & 40 & catalyseur (0,7\% en masse) \\
\hline
\end{tabularx}
\end{center}

\end{document}
